\naslov{Preizkušanje domnev}

Želimo narediti nek sklep o $\theta$, včasih imamo o njem neko idejo.
Tedaj to idejo imenujemo \pojem{ničelna hipoteza}, jo označimo s $H_0$ in
poskusimo preveriti ali zavreči.
V splošnem $H_0$ trdi, da je $\theta \in I_0 \subseteq I$.
Testiramo jo proti \pojem{alternativni domnevi} $H_1$.
Če je opažanje preveč neskladno s $H_0$, to domnevo zavrnemo in sprejmemo $H_1$;
nikoli pa ne sprejmemo $H_0$.

V primeru, ko $H_0$ zavrnemo, čeprav velja, pride do zmote.
Želimo si, da je takih zmot malo, hkrati pa želimo $H_0$ čim večkrat zavreči, če
ne velja.
Funkciji $I_1 \to [0,1]$ s predpisom
\[
  \theta \mapsto P_\theta(\text{zavrnemo}\,H_0)
\]
pravimo \pojem{moč preizkusa}, in nam pove, do kolikšne mere smo upoštevali
drugo od teh želja.
Za prvo željo običajno le želimo omejiti verjetnost zmote
$P_\theta(\text{zavrnemo}\,H_0)$ na $\alpha$ za $\theta \in I_0$.
Tu uporabimo podobno idejo kot pri intervalih zaupanja.

\vprasanje{Opiši osnovno idejo za testiranjem hipotez. Kaj je moč preizkusa?}

Včasih pride prav pogledati razmerje verjetij
\[
  \Lambda = \frac{\sup_{\theta \in I} L(\theta \such X)}{\sup_{\theta \in I_0}
	L(\theta \such X)},
\]
kjer bomo hipotezo zavrnili, če bo razmerje dovolj veliko.
Ker je porazdelitev tega v praksi zelo težko poiskati, si pomagamo z naslednjim
izrekom:

\begin{izrek}[Wilks]
  Naj bodo $X_1, X_2, \ldots$ neodvisne in enako porazdeljene, njihova
  porazdelitev pa naj pripada družini $P_\theta$, $\theta \in I$, kjer je $I$
  gladka mnogoterost dimenzije $p$.
  Naj bo $I_0$ gladka podmnogoterost dimenzije $q$.
  Tedaj velja
  \[
	2 \log \Lambda \xrightarrow[n \to \infty]{d} \chi^2(p-q)
  \]
  pri vsaki verjetnostni meri $P_\theta$ za $\theta \in I_0$.
\end{izrek}

\vprasanje{Kaj je razmerje verjetij? Povej Wilksov izrek.}

% LocalWords:  Wilksov
